\chapter{Allgemeines zur Virtualisierung}
\label{kap:allgemeines}

\begin{aufgabe}[Was ist eigentlich ein Hypervisor?]
Der Hypervisor ist so etwas wie ein Magier, der virtuelle Maschinen zum Leben erweckt. Aber es gibt zwei verschiedene Typen von Hypervisoren! Erkläre, was ein Hypervisor ist und warum VirtualBox ein \enquote{Typ-2-Magier} ist.
\end{aufgabe}

% TODO Antwort

\begin{aufgabe}[Was ist das Host-System und was das Guest-System?]
Wenn du VirtualBox nutzt, arbeitest du sowohl mit deinem normalen Betriebssystem (Host-System) als auch mit einem \enquote{virtuellen} Betriebssystem (Guest-System). Erkläre den Unterschied zwischen dem Host- und dem Guest-System und wie sie miteinander interagieren.
\end{aufgabe}

% TODO Antwort

\begin{aufgabe}[Was sind Guest Additions, und warum sind sie nützlich?]
Mit den Guest Additions bekommt deine virtuelle Maschine Superkräfte! Finde heraus, was Guest Additions sind und warum du sie installieren solltest. Nenne 6 Punkte bzw. Aufgaben, welche durch die Guest-Additions erfüllt werden.
\end{aufgabe}

% TODO Antwort

\begin{aufgabe}[Warum müssen Guest-Additions überhaupt zusätzlich installiert werden?]
Guest-Additions müssen im Guest-System zusätzlich installiert werden. Begründe warum das notwendig ist und diese Aufgaben nicht direkt von VirtualBox durchgeführt werden können.
\end{aufgabe}

% TODO Antwort

\begin{aufgabe}[Was ist der Host-Key und wozu wird er verwendet?]
In VirtualBox gibt es eine spezielle Taste mit der du besondere Funktionen auslösen kannst. Was ist das für eine Taste? Nenne weiters mindestens 4 Funktionen wozu du diese benötigst.
\end{aufgabe}

% TODO Antwort

\begin{aufgabe}[Was ist eine virtuelle Festplatte, und welche Arten gibt es?]
Eine virtuelle Maschine braucht auch eine \enquote{virtuelle Festplatte}. Welche Arten von virtuellen Festplatten gibt es und was sind die Unterschiede?
\end{aufgabe}

% TODO Antwort

\begin{aufgabe}[Was ist ein Shared Folder in VirtualBox und wozu wird er verwendet?]
\end{aufgabe}

% TODO Antwort

\begin{aufgabe}[Was sind Snapshots, und wofür werden sie verwendet?]
Stell dir vor, du könntest deine virtuelle Maschine \enquote{einfrieren} und später genau an dieser Stelle weitermachen. Genau das machen Snapshots! Erkläre, was Snapshots in VirtualBox sind und wie sie dir helfen können, Fehler zu vermeiden.
\end{aufgabe}

% TODO Antwort

\begin{aufgabe}[Warum glaubst du, ist ein Snapshot für uns im Labor besonders wichtig?]
\end{aufgabe}

% TODO Antwort
